% Source: tex/chapters/ch05/07_テスト関連測度.tex
\subsubsection{実施したテストの評価}
テストそのものの評価では、テストスイートがどれだけ強いか、どの技法が有効かを比較する。カバレッジ、ミューテーション分析、欠陥注入などが用いられる。

\paragraph{欠陥注入}
欠陥注入は、人工的な欠陥を SUT に入れ、テストスイートがそれらをどれだけ検出できるかを調べる。検出率からテスト有効性や残存欠陥数を推定しようとする。ただし、人工欠陥が本物の欠陥を代表しているか、挿入した欠陥を取り残さないかには注意が必要である。

\paragraph{ミューテーションスコア}
ミューテーションスコアは、殺された変異体数を生成した変異体数で割った値である。値が高いほど、テストスイートが小さな変化に敏感であることを示す。ただし、変異体の生成数が多くなるため、実行コストが問題になる。

\paragraph{技法間の比較と相対有効性}
相対有効性は、異なるテスト技法を同じ性質に対して比較する考え方である。最初の障害を見つけるまでのテスト数、欠陥検出率、信頼性改善量などが比較軸になる。
